\chapter{Wymagania i koncepcja systemu}\label{ch:wymagania}
% Ok. 5-7 stron. Materiał: HD 'Założenia architektoniczne'.
\begin{zawartosc}
  \item Lista wymagań funkcjonalnych i niefunkcjonalnych w tabeli z numerami (weryfikowane w rozdz. 7).
  \item Architektura systemu na schemacie blokowym: płytka główna (MB), daughterboard z USB-C (DB), dwie folie membrany, ogniwo Li-Po.
  \item Wybór kluczowych układów (nRF52840, nPM1300) z porównaniem alternatyw i uzasadnieniem: bare chip zamiast modułu, PMIC zamiast osobnej ładowarki i regulatora.
  \item Wstępne bilanse: liczba GPIO, budżet prądowy szyny 3,3 V, szacowany czas pracy z danych katalogowych.
\end{zawartosc}

\section{Wymagania funkcjonalne oraz niefunkcjonalne}\label{sec:wymagania}
% - Tabela W1..Wn, każde wymaganie weryfikowane w sekcji 7.7.
% - Funkcjonalne: BLE + USB HID, NKRO, ładowanie z USB-C, dioda statusu, SWD, raport baterii.
% - Niefunkcjonalne: drop-in, czas pracy (wartość do ustalenia), produkcja JLC, koszt, rozszerzalność (J3).
\todo{Napisać.}

\section{Architektura systemu}\label{sec:architektura}
% - Schemat blokowy: MB, DB, łącze 1x10, membrana (J4, J7), ogniwo (J2). Główny rysunek pracy.
% - Podział na płytki i projekty KiCada (CL A2).
\todo{Napisać.}

\section{Wybór układów}\label{sec:wybor}
% - MCU: nRF52840 bare chip vs moduł; tabela alternatyw.
% - PMIC: nPM1300 vs ładowarka + LDO/buck; tabela alternatyw.
% - Ogniwo i złącze (CL A26, A27, P11).
% - Budżet: GPIO (HD 'Przydział GPIO'), prąd szyny 3,3 V (HD 'Obciążenie szyny 3,3 V'), szacunek czasu pracy.
\todo{Napisać.}

